\documentclass[11pt]{article}
\usepackage{amssymb}
\usepackage{amsthm}
\usepackage{enumitem}
\usepackage{physics,amsmath}
\usepackage{bm}
\usepackage{adjustbox}
\usepackage{mathrsfs}
\usepackage{graphicx}
\usepackage{siunitx}
\usepackage[mathscr]{euscript}


\title{\textbf{Solved selected problems of General Relativity - Thomas A. Moore}}
\author{Franco Zacco}
\date{}

\addtolength{\topmargin}{-3cm}
\addtolength{\textheight}{3cm}

\newcommand{\hatr}{\bm{\hat{r}}}
\newcommand{\hatn}{\bm{\hat{n}}}
\newcommand{\hatx}{\bm{\hat{x}}}
\newcommand{\haty}{\bm{\hat{y}}}
\newcommand{\hatz}{\bm{\hat{z}}}
\newcommand{\hatth}{\bm{\hat{\theta}}}
\newcommand{\hatphi}{\bm{\hat{\varphi}}}
\newcommand{\hatrho}{\bm{\hat{\rho}}}
\newcommand{\esi}[1]{\bm{e}_{#1}}
\newcommand{\etht}{\bm{e}_\theta}
\newcommand{\sci}[1]{\times 10^{#1}}

\theoremstyle{definition}
\newtheorem*{solution*}{Solution}
\renewcommand*{\proofname}{Solution}

\begin{document}
\maketitle
\thispagestyle{empty}

\section*{Chapter 20 - The Stress-Energy Tensor}

\begin{proof}{\textbf{BOX 20.1} - Exercise 20.1.1.}\\
If the source of the gravitational field of the box were the total energy,
then we would not be firing particles to infinity with energy we create out of
nothing, since after the matter-antimatter combination, the photons created
would hold kinetic energy, which would still be in the box because it's
perfectly mirrored, in short, because of the conservation of energy, we cannot
create energy out of nothing.
\end{proof}

\begin{proof}{\textbf{BOX 20.2} - Exercise 20.2.1.}\\
Let us consider the Stress-Energy components $T^{ij} = \rho_0 u^i u^j$ where
$i$ and $j$ are both spatial indices, then we see that
\begin{align*}
    T^{ij} &= \rho_0 u^i u^j = n_0 m u^i u^j 
    = n_0 m (u^tv_i) u^j
    = (n_0 u^t)v_i (m u^j) =\\
    &= n v_i p^j
    = \frac{(n A v_i dt) p^j}{A dt}
\end{align*}
Therefore $T^{ij}$ is the i-flux of j-momentum.
\end{proof}

\cleardoublepage
\begin{proof}{\textbf{BOX 20.3} - Exercise 20.3.1.}\\
If we consider the component $T^{xx}$ then the integral over $p^x$ becomes
\begin{align*}
    \int_{-\infty}^\infty f([p^x]^2) (p^x)^2 ~dp^x
\end{align*}
We note that the integrand is an even function, so the integral along a
symmetric interval is not necessarily 0, as opposed to the case where the 
integrand is an odd function like $f([p^x]^2) p^x$.
\end{proof}

\begin{proof}{\textbf{BOX 20.3} - Exercise 20.3.2.}\\
Since the selection of the $x$ direction was arbitrary must be true for any
other direction. Hence be symmetry we have that $T^{xx} = T^{yy} = T^{zz}$.
\end{proof}

\begin{proof}{\textbf{BOX 20.3} - Exercise 20.3.3.}\\
We note that $N(|\vec{p}|)/L^3$ is the number density of particles having
momentum $\vec{p}$ and $p^t$ is the energy of a single particle with momentum
$\vec{p}$, then $N(|\vec{p}|)p^t/L^3$ is the energy density of particles with
momentum $\vec{p}$.
\\
Then, in the integral
\begin{align*}
    T^{tt} = \frac{1}{L^3}\iiint dp^xdp^ydp^z N(|\vec{p}|)p^t
\end{align*}
We are summing across the whole momentum space i.e. we are considering all
possible momentums inside the box so we get that $T^{tt} = \rho$ where $\rho$
is the total energy density inside the box.
\end{proof}

\cleardoublepage
\begin{proof}{\textbf{BOX 20.4} - Exercise 20.4.1.}\\
In a LIF where the fluid at a given event is (at least instantaneously) at rest
the metric $g_{\mu\nu} = \eta_{\mu\nu}$ and the fluid's four-velocity at that
event is $u^\mu = [1,0,0,0]$.\\
Given that $T^{\mu\nu} = (\rho_0 + p_0)u^\mu u^\nu + p_0g^{\mu\nu}$ then if
$\mu \neq \nu$ then $T^{\mu\nu} = 0$ because the first term vanishes and the
only non-zero components of $\eta^{\mu\nu}$ are the ones where $\mu = \nu$.\\
If $\mu = \nu = t$ we see that
\begin{align*}
    T^{tt} = (\rho_0 + p_0)u^t u^t + p_0\eta^{tt}
    = (\rho_0 + p_0)\cdot 1 + p_0\cdot (-1)
    = \rho_0
\end{align*}
And if $\mu = \nu = x \text{ or } y \text{ or } z$ we get that 
\begin{align*}
    T^{xx} = (\rho_0 + p_0)u^x u^x + p_0\eta^{xx} = 0 + p_0 \cdot 1
    = p_0\\
    T^{yy} = (\rho_0 + p_0)u^y u^y + p_0\eta^{yy} = 0 + p_0 \cdot 1
    = p_0\\
    T^{zz} = (\rho_0 + p_0)u^z u^z + p_0\eta^{zz} = 0 + p_0 \cdot 1
    = p_0
\end{align*}
Therefore 
\begin{align*}
    T^{\mu\nu} = \begin{bmatrix}
        \rho_0 & 0 & 0 & 0\\
        0 & p_0 & 0 & 0\\
        0 & 0 & p_0 & 0\\
        0 & 0 & 0 & p_0
    \end{bmatrix}
\end{align*}
\end{proof}

\cleardoublepage
\begin{proof}{\textbf{BOX 20.5} - Exercise 20.5.1.}\\
Equation (20.29) states that
\begin{align*}
    u^{\mu}u^{\nu}\partial_{\mu}(\rho_0 + p_0)
    + (\rho_0 + p_0)[u^{\nu}\partial_{\mu}u^\mu + u^\mu \partial_\mu u^\nu] 
    + \eta^{\mu\nu}\partial_\mu p_0 = 0
\end{align*}
So multiplying the equation by $u_\nu$ we get that
\begin{align*}
    u_{\nu}u^{\mu}u^{\nu}\partial_{\mu}(\rho_0 + p_0)
    + (\rho_0 + p_0)[u_{\nu}u^{\nu}\partial_{\mu}u^\mu + u_{\nu}u^\mu \partial_\mu u^\nu] 
    + u_{\nu}\eta^{\mu\nu}\partial_\mu p_0
    &= 0\\
    -u^{\mu}\partial_{\mu}(\rho_0 + p_0)
    + (\rho_0 + p_0)[-\partial_{\mu}u^\mu + u_{\nu}u^\mu \partial_\mu u^\nu] 
    + u^{\mu}\partial_\mu p_0
    &= 0
\end{align*}
Where we used that $u^\nu u_\nu = -1$.\\
Now, using that $u_\nu\partial_{\mu}u^\nu = 0$ we see that
\begin{align*}
    -u^{\mu}\partial_{\mu}(\rho_0 + p_0)
    + (\rho_0 + p_0)[-\partial_{\mu}u^\mu + 0] 
    + u^{\mu}\partial_\mu p_0
    &= 0\\
    -u^{\mu}\partial_{\mu}\rho_0 - u^{\mu}\partial_{\mu}p_0
    -\rho_0\partial_{\mu} u^\mu - p_0\partial_{\mu}u^\mu
    + u^{\mu}\partial_\mu p_0
    &= 0\\
    -u^{\mu}\partial_{\mu}\rho_0
    -\rho_0\partial_{\mu} u^\mu - p_0\partial_{\mu}u^\mu
    &= 0\\
    \partial_{\mu}(\rho_0 u^{\mu}) + p_0\partial_{\mu}u^\mu
    &= 0
\end{align*}
\end{proof}

\begin{proof}{\textbf{BOX 20.5} - Exercise 20.5.2.}\\
Multiplying the equation $\partial_{\mu}(\rho_0 u^{\mu}) + p_0\partial_{\mu}u^\mu = 0$
by $u^\nu$ we get that
\begin{align*}
    u^\nu\partial_{\mu}(\rho_0 u^{\mu}) + u^\nu p_0\partial_{\mu}u^\mu
    &= 0
\end{align*}
Subtracting this equation from equation (20.29) we get that
\begin{align*}
    &u^{\mu}u^{\nu}\partial_{\mu}(\rho_0 + p_0)
    + (\rho_0 + p_0)[u^{\nu}\partial_{\mu}u^\mu + u^\mu \partial_\mu u^\nu] 
    + \eta^{\mu\nu}\partial_\mu p_0\\
    &- u^\nu u^{\mu}\partial_{\mu}\rho_0 - u^\nu \rho_0\partial_{\mu}u^\mu 
    - u^\nu p_0\partial_{\mu}u^\mu
    = 0\\[6pt]
    &u^{\mu}u^{\nu}\partial_{\mu}p_0
    + (\rho_0 + p_0)u^\mu \partial_\mu u^\nu
    + \eta^{\mu\nu}\partial_\mu p_0
    = 0\\[6pt]
    &(u^{\mu}u^{\nu} + \eta^{\mu\nu})\partial_{\mu}p_0
    + (\rho_0 + p_0)u^\mu \partial_\mu u^\nu
    = 0\\
\end{align*}
Therefore
\begin{align*}
    (\rho_0 + p_0)u^\mu \partial_\mu u^\nu
    = -(u^{\mu}u^{\nu} + \eta^{\mu\nu})\partial_{\mu}p_0
\end{align*}
\end{proof}

\cleardoublepage
\begin{proof}{\textbf{P20.1}}\\
Let us consider the ionized hydrogen plasma at the sun's center, where the
density is $160~g/cm^3$ and the temperature is $1.5\sci{7}~K$.
\\
We can compute the energy of an individual electron in the plasma by
\begin{align*}
    p^t_e &= m_e + \frac{3}{2}k_B T\\
    &= 9.109\sci{-31}~kg + \frac{3}{2}\cdot 1.536\sci{-49}~kg/K \cdot 1.5\sci{7}~K\\
    &\approx 9.109\sci{-31}~kg
\end{align*}
In the same way for a proton we ge that $p^t_p \approx m_p = 1.672\sci{-27}{kg}$.
\\
Also, we can compute the number density of particles as follows
\begin{align*}
    n = 160~\frac{g}{cm^3} \frac{2 \cdot 6.022\sci{23}~particles}{1~g}
    = 1.927\sci{26}~\frac{particles}{cm^3}
\end{align*}
Since the particles are not moving relativistically then we may assume $u^t = 1$
and hence $\rho = n p^t = \rho_0$ so the energy density is
\begin{align*}
    \rho_0 &= \frac{n}{2} p^t_e + \frac{n}{2} p^t_p\\
    &= \frac{1.927\sci{26}}{2} \cdot 9.109\sci{-31}
    + \frac{1.927\sci{26}}{2} \cdot 1.672\sci{-27}\\
    &= 8.776\sci{-5} + 0.161\\
    &= 0.16108~kg/cm^3 = 161.08 ~g/cm^3
\end{align*}
We note that $\rho_0 \approx \text{mass density}$.
\\
Now, using the ideal gas law we see that
\begin{align*}
    p_0 &= \frac{N}{V} k_B T \\
    &= n k_B T\\
    &= 1.927\sci{26} \cdot 1.536\sci{-46} \cdot 1.5\sci{7}\\
    &= 4.439\sci{-13}~g/cm^3
\end{align*}
Therefore we note that $p_0 \ll \rho_0$ and we have all the components of the 
stress-energy tensor.
\end{proof}

\cleardoublepage
\begin{proof}{\textbf{P20.2}}\\
From equation (20.21) we know that $p_0 = T^{xx} = \rho_0 (u^{x})^2$
we want that $\rho_0 = 100 p_0$ so we get that
\begin{align*}
    p_0 &= 100 p_0 (u^{x})^2\\
    u^{x} &= \sqrt{\frac{1}{100}} = 0.1
\end{align*}
In the same way, we see that $u^{y} = u^{z} = 0.1$. Then we get that
\begin{align*}
    v_x = u^x\sqrt{1 - v^2} = 0.1\sqrt{1 - v^2}
\end{align*}
In the same way we have that $v_y = v_z = 0.1\sqrt{1 - v^2}$.
Then we solve for $v$ as follows 
\begin{align*}
    v &= \sqrt{v_x^2 + v_y^2 + v_z^2}\\
    v &= \sqrt{3\cdot 0.01(1 - v^2)}\\
    v^2 &= 0.03 - 0.03v^2\\
    v^2(1 + 0.03) &= 0.03\\
    v &= \sqrt{\frac{0.03}{1.03}} = 0.171
\end{align*}
\end{proof}

\cleardoublepage
\begin{proof}{\textbf{P20.3}}\\
Given that the star's pressure is negligible then equation (20.16) becomes
\begin{align*}
    T^{\mu\nu} = \rho u^\mu u^\nu
\end{align*}
Also, the velocity at any point of the star, in spherical coordinates has 
the form
\begin{align*}
    \bm{v} = r\omega \sin\theta \hatphi
\end{align*}
or in cartesian coordinates (the LIF coordinates)
\begin{align*}
    \bm{v} &= r\omega \sin\theta(-\sin\varphi \hatx + \cos\varphi \haty)\\
    &= -y\omega\hatx + x\omega\haty
\end{align*}
Where we used that $x = r\sin\theta\cos\varphi$ and $y = r\sin\theta\sin\varphi$.
\\
Then, the four-velocity $u^\alpha$ of the star, ignoring the terms of order $v^2$,
gives us
\begin{align*}
    u^\alpha = \begin{bmatrix}
        1 \\ v_x \\ v_y \\ 0
    \end{bmatrix}
    =
    \begin{bmatrix}
        1 \\ -y\omega \\ x\omega \\ 0
    \end{bmatrix}
\end{align*}
Therefore the stress-energy tensor for the rotating star is
\begin{align*}
    T^{\mu\nu} = \begin{bmatrix}
        \rho & -\rho y\omega & \rho x\omega & 0\\
        -\rho y\omega & \rho y^2\omega^2 & -\rho yx\omega^2 & 0\\
        \rho x\omega & -\rho yx\omega^2 & \rho x^2\omega^2 & 0\\
        0 &  0 & 0 & 0
    \end{bmatrix}
\end{align*}
But this form of the stress-energy tensor still has components of order $v_x^2$
and $v_y^2$, so dropping them we get that
\begin{align*}
    T^{\mu\nu} = \begin{bmatrix}
        \rho & -\rho y\omega & \rho x\omega & 0\\
        -\rho y\omega & 0 & 0 & 0\\
        \rho x\omega & 0 & 0 & 0\\
        0 &  0 & 0 & 0
    \end{bmatrix}
\end{align*}
\end{proof}

\cleardoublepage
\begin{proof}{\textbf{P20.4}}\\
Let us consider an ideal gas of photons in a perfectly-mirrored box in a LIF
at rest with respect to the box.\\
We know that
\begin{align*}
    T^{\mu\nu} = \frac{1}{L^3}\iiint dp^x dp^y dp^z N(|\vec{p}|) \frac{p^\mu p^\nu}{p^t}
\end{align*}
In particular if $\mu = \nu = t$ then we have that
\begin{align*}
    T^{tt} = \frac{1}{L^3}\iiint dp^x dp^y dp^z N(|\vec{p}|)p^t
\end{align*}
We know that $N(|\vec{p}|)/L^3$ is the number density of photons having
momentum $\vec{p}$ and $p^t = E$ is the energy of a single photon with momentum
$\vec{p}$, then $N(|\vec{p}|)p^t/L^3$ is the energy density of particles with
momentum $\vec{p}$.\\
Then the integral is the sum of all energy densities across the whole momentum
space so $T^{tt} = \rho$ where $\rho$ is the total energy density inside the box.
\\
Now, if $\mu = \nu = x$ we have that
\begin{align*}
    T^{xx}
    &= \frac{1}{L^3}\iiint dp^x dp^y dp^z N(|\vec{p}|) \frac{p^x p^x}{p^t}\\
    &= \frac{1}{L^3}\iiint dp^x dp^y dp^z N(|\vec{p}|) \frac{E^2v_x^2}{E}\\
    &= \frac{1}{L^3}\iiint dp^x dp^y dp^z N(|\vec{p}|) E v_x^2
\end{align*}
Where we used that $p^x = Ev_x$, in the same way we see that
\begin{align*}
    T^{yy} &= \frac{1}{L^3}\iiint dp^x dp^y dp^z N(|\vec{p}|) E v_y^2\\
    T^{zz} &= \frac{1}{L^3}\iiint dp^x dp^y dp^z N(|\vec{p}|) E v_z^2
\end{align*}
Hence
\begin{align*}
    T^{xx} + T^{yy} + T^{zz}
    &= \frac{1}{L^3}\iiint dp^x dp^y dp^z N(|\vec{p}|) E (v_x^2 + v_y^2 + v_z^2)\\
    &= \frac{1}{L^3}\iiint dp^x dp^y dp^z N(|\vec{p}|) E\\
    &= \rho
\end{align*}
Where we used that $v_x^2 + v_y^2 + v_z^2 = 1$ for photons.\\
Also, we know that for an ideal gas $T^{xx} = T^{yy} = T^{zz}$, so we get that
\begin{align*}
    T^{xx} = \frac{1}{3}\rho
    \qquad
    T^{yy} = \frac{1}{3}\rho
    \qquad
    T^{zz} = \frac{1}{3}\rho
\end{align*}
We see that the same mathematical argument works to show that all off-diagonal
elements in the stress-energy tensor are 0.\\
Therefore the stress-energy tensor for an ideal gas of photons is
\begin{align*}
    T^{\mu\nu} = \begin{bmatrix}
        \rho & 0 & 0 & 0\\
        0 & \frac{1}{3} \rho & 0 & 0\\
        0 & 0 & \frac{1}{3} \rho & 0\\
        0 & 0 & 0 & \frac{1}{3} \rho
    \end{bmatrix}
\end{align*}
\end{proof}

\cleardoublepage
\begin{proof}{\textbf{P20.5}}\\
\begin{itemize}
\item [\textbf{a.}] We know that
\begin{align*}
    F^{\mu\nu}
    = \begin{bmatrix}
        0 & E_x & E_y & E_z\\
        -E_x & 0 & B_z & -B_y\\ 
        -E_y & -B_z & 0 & B_x\\ 
        -E_z & B_y & -B_x & 0\\ 
    \end{bmatrix}
\end{align*}
Before computing $F^{\mu\nu}F_{\mu\nu}$ we need to note 
first that in a LIF $F_{\mu\nu}$ is given by
\begin{align*}
    F_{\alpha\beta} = \eta_{\alpha\mu} \eta_{\beta\nu}F^{\mu\nu}
    = \begin{bmatrix}
        0 & -E_x & -E_y & -E_z\\
        E_x & 0 & B_z & -B_y\\ 
        E_y & -B_z & 0 & B_x\\ 
        E_z & B_y & -B_x & 0\\ 
    \end{bmatrix}
\end{align*}
Now, we see that
\begin{align*}
    F^{\mu\nu}F_{\mu\nu}
    &= F^{tx}F_{tx} + F^{ty}F_{ty} + F^{tz}F_{tz}
    + F^{xt}F_{xt} + F^{xy}F_{xy} + F^{xz}F_{xz}\\
    &\quad + F^{yt}F_{yt} + F^{yx}F_{yx} + F^{yz}F_{yz}
    + F^{zt}F_{zt} + F^{zx}F_{zx} + F^{zy}F_{zy}\\
    %
    &= -E_x^2 -E_y^2  -E_z^2
    -E_x^2 + B_z^2 + B_y^2
    -E_y^2 + B_z^2 + B_x^2
    -E_z^2 + B_y^2 + B_x^2\\
    &= 2(B^2 - E^2)
\end{align*}
Where we used that $E^2 = E_x^2 + E_y^2 + E_z^2$ and
$B^2 = B_x^2 + B_y^2 + B_z^2$.\\
Let us consider now a second-rank tensor defined as 
\begin{align*}
    T^{\mu\nu} = F^{\mu\alpha}F^{\nu}_{\alpha}
    + Cg^{\mu\nu}F^{\alpha\beta}F_{\alpha\beta}
\end{align*}
We see that
\begin{align*}
    F^{\nu}_{\alpha} &= \eta_{\alpha\mu} F^{\mu\nu}
    = \begin{bmatrix}
        0 & -E_x & -E_y & -E_z\\
        -E_x & 0 & B_z & -B_y\\ 
        -E_y & -B_z & 0 & B_x\\ 
        -E_z & B_y & -B_x & 0\\ 
    \end{bmatrix}
\end{align*}
Then if we compute $T^{tt}$ in a LIF we get that
\begin{align*}
    T^{tt}
    &= F^{t\alpha}F^{t}_{\alpha} + C\eta^{tt}F^{\alpha\beta}F_{\alpha\beta}\\
    &= -E_x^2 - E_y^2 - E_z^2  -2C(B^2 - E^2)\\
    &= -E^2 -2C(B^2 - E^2)\\
    &= -2CB^2 + E^2(2C - 1)
\end{align*}
We need that $-2C = 2C - 1$ so $C = 1/4$ and hence must be that
\begin{align*}
    T^{\mu\nu} = -\frac{1}{4\pi k} \bigg(F^{\mu\alpha}F^{\nu}_{\alpha}
    + \frac{1}{4}g^{\mu\nu}F^{\alpha\beta}F_{\alpha\beta}\bigg)
\end{align*}
So when $\mu = \nu = t$ we get that
\begin{align*}
    T^{tt} = -\frac{1}{4\pi k} \bigg(-\frac{1}{2}(E^2 + B^2)\bigg)
    = \frac{1}{8\pi k}(E^2 + B^2) = \rho_E
\end{align*}
as we wanted.\\
Also, we need to check if the defintion of $T^{\mu\nu}$ we gave is symmetric,
since the metric $g^{\mu\nu}$ is symmetric we only need to show that
$F^{\mu\alpha}F^{\nu}_{\alpha}$ is symmetric.\\
We see that
\begin{align*}
    F^{t\alpha}{F_{\alpha}}^{t} &= -E_x^2 -E_y^2 - E_z^2 = -E^2\\
    F^{t\alpha}{F_{\alpha}}^{x} &= -E_yB_z + E_zB_y\\
    F^{t\alpha}{F_{\alpha}}^{y} &= E_xB_z - E_zB_x\\
    F^{t\alpha}{F_{\alpha}}^{z} &= -E_xB_y + E_yB_x
\end{align*}
\begin{align*}
    F^{x\alpha}{F_{\alpha}}^{t} &= -B_zE_y + B_yE_z = F^{t\alpha}{F_\alpha}^{x}\\
    F^{x\alpha}{F_{\alpha}}^{x} &= E_x^2 - B_z^2 - B_y^2\\
    F^{x\alpha}{F_{\alpha}}^{y} &= E_xE_y + B_xB_y\\
    F^{x\alpha}{F_{\alpha}}^{z} &= E_xE_z + B_xB_z
\end{align*}
\begin{align*}
    F^{y\alpha}{F_{\alpha}}^{t} &= E_xB_z - E_zB_x = F^{t\alpha}{F_{\alpha}}^{y}\\
    F^{y\alpha}{F_{\alpha}}^{x} &= E_yE_x + B_xB_z = F^{x\alpha}{F_{\alpha}}^{y}\\
    F^{y\alpha}{F_{\alpha}}^{y} &= E_y^2 - B_z^2 - B_x^2\\
    F^{y\alpha}{F_{\alpha}}^{z} &= E_yE_z + B_yB_z
\end{align*}
\begin{align*}
    F^{z\alpha}{F_{\alpha}}^{t} &= -E_xB_y + E_yB_x = F^{t\alpha}{F_{\alpha}}^{z}\\
    F^{z\alpha}{F_{\alpha}}^{x} &= E_xE_z + B_xB_z = F^{x\alpha}{F_{\alpha}}^{z}\\
    F^{z\alpha}{F_{\alpha}}^{y} &= E_yE_z + B_yB_z = F^{y\alpha}{F_{\alpha}}^{z}\\
    F^{z\alpha}{F_{\alpha}}^{z} &= E_z^2 - B_y^2 - B_x^2
\end{align*}
Therefore since $F^{\mu\alpha}F^{\nu}_{\alpha}$ is symmetric then $T^{\mu\nu}$
is symmetric.
\item [\textbf{b.}] Let us compute $T^{tx}$ as follows
\begin{align*}
    T^{tx} &= - \frac{1}{4\pi k} \bigg(
        F^{t\alpha}{F_{\alpha}}^{x}
        + \frac{1}{4}\eta^{tx}F^{\alpha\beta}F_{\alpha\beta}
    \bigg)\\
    &= - \frac{1}{4\pi k} \bigg(E_zB_y - E_yB_z\bigg)\\
    &= - \frac{1}{4\pi k} \bigg(\vec{E} \times \vec{B}\bigg)_x
\end{align*} 
\end{itemize}
\end{proof}

\cleardoublepage
\begin{proof}{\textbf{P20.6}}
\begin{itemize}
\item [\textbf{a.}] We know that if $i$ and $j$ are spatial coordinates
then $T^{ij}$ is the $i$-flux of $j$-momentum.\\
On the other hand, we saw that if a plate with area $A$ is placed perpendicular
to the $x$ direction (in the path of the dust particles) then, for example,
$T^{yx}A$ is the $y$-momentum per unit of coordinate time $t$.\\\\
Suppose we are in a LIF.
If the plate is in any position we can generalize this by taking a unit normal
four-vector perpendicular to $A$ with components $\bm{n} = [0, n_x, n_y, n_z]$
where the time component is 0 since it's purely a spatial vector.\\
Then $T^{\mu\nu} A n_\nu$ is the $\mu$-momentum per unit of time, where $\mu$
is a spatial coordinate.\\
If $\mu$ is a time coordinate then $T^{t\nu}$ is the $\nu$-flux of energy so
$T^{t\nu} A n_\nu$ is the $\nu$-energy per unit of time.
Therefore the four-force $\bm{F}$ is given by
\begin{align*}
    \bm{F} = \bigg(\dv{\bm{p}}{\tau}\bigg)^\mu = T^{\mu\nu} A n_\nu
    = T^{\mu\nu} A \eta_{\nu\alpha} n^\alpha
\end{align*}
Where we added in the last step $\eta_{\nu\alpha}$ since we are working on a
LIF and it doesn't change the equation since the time component of $\bm{n}$ is
0.\\
Finally, we can generalize this, using the appropiate metric as 
\begin{align*}
    \bm{F} = \bigg(\dv{\bm{p}}{\tau}\bigg)^\mu
    = T^{\mu\nu} A g_{\nu\alpha} n^\alpha
\end{align*}

\item [\textbf{b.}] Let us consider a perfect fluid. If $\bm{n}$ represents
a wall in space and $\bm{u}$ represents the fluid's four velocity then if
the fluid is at rest with respect to the wall we must have that
$\bm{u} \cdot \bm{n} = g_{\nu\alpha}u^{\nu} n^\alpha = 0$ since the only
non-zero component of $u^\nu$ is $u^t = 1$ but $n^t = 0$.\\
Then the four-force exerted on any patch of area $A$ on the wall of the
container holding the perfect fluid at rest is
\begin{align*}
    \bm{F} = \bigg(\dv{\bm{p}}{\tau}\bigg)^\mu
    &= T^{\mu\nu} A g_{\nu\alpha} n^\alpha\\
    &= [(\rho_0 + p_0)u^\mu u^\nu + p_0 g^{\mu\nu}]A g_{\nu\alpha} n^\alpha\\   
    &= 0 + p_0 g^{\mu\nu}A g_{\nu\alpha} n^\alpha\\
    &= p_0 A \delta^{\mu}_{\alpha} n^\alpha\\
    &= p_0 A n^\mu
\end{align*}
Therefore the magnitude of the four-force exerted on any patch of area $A$
on the wall of the container is $p_0 A$.
\end{itemize}
\end{proof}

\cleardoublepage
\begin{proof}{\textbf{P20.7}}\\
Let us consider the tensor $T^{\mu\nu} = - \Lambda g^{\mu\nu}$ if we take a LIF
then the tensor becomes $T^{\mu\nu} = - \Lambda \eta^{\mu\nu}$ i.e.
\begin{align*}
    T^{\mu\nu} = \begin{bmatrix}
        \Lambda & 0 & 0 & 0\\
        0 & -\Lambda & 0 & 0\\
        0 & 0 & -\Lambda & 0\\
        0 & 0 & 0 & -\Lambda
    \end{bmatrix}
\end{align*}
Comparing this Stress-Energy tensor to the one we found for a perfect fluid
then we see that the pressure $p_0 = -\Lambda$ and since lambda is a positive
scalar then the pressure is negative, so if $T^{\mu\nu}$ is a Stress-Energy 
tensor of a real fluid must be, for example, that
\begin{align*}
    T^{xx} = -\Lambda
    = \frac{1}{L^3}\iiint N(|\vec{p}|)\frac{(p^x)^2}{p^t}~dp^x dp^y dp^z
\end{align*}
But $N(|\vec{p}|) \geq 0$, $(p^x)^2 \geq 0$ and $p^t = m/\sqrt{1 - v^2}$ which
is also positive, so the integral must be positive as well.\\
Therefore $T^{\mu\nu} = - \Lambda g^{\mu\nu}$ cannot represent a real fluid's 
Stress-Energy tensor.
\end{proof}

\cleardoublepage
\begin{proof}{\textbf{P20.8}}
\begin{itemize}
\item [\textbf{a.}]
Let us consider the second-rank tensor $M_{\alpha\beta,obs}$
defined as the tensor product of two four-vectors
\begin{align*}
    M_{\alpha\beta,obs} &= A_{\alpha, obs} B_{\beta, obs}\\
    &= (\bm{o}_{\alpha})^\sigma g_{\sigma\gamma} A^{\gamma}
    (\bm{o}_{\beta})^\lambda g_{\lambda\delta} B^{\delta}
\end{align*}
We also note that
\begin{align*}
    M^{\mu\nu}_{obs}
    &= \eta^{\mu\alpha} \eta^{\nu\beta} M_{\alpha\beta,obs} \\
    &= \eta^{\mu\alpha} \eta^{\nu\beta}
    (\bm{o}_{\alpha})^\sigma g_{\sigma\gamma} A^{\gamma}
    (\bm{o}_{\beta})^\lambda g_{\lambda\delta} B^{\delta} \\
    &= \eta^{\mu\alpha} \eta^{\nu\beta}A^{\gamma}B^{\delta}
    g_{\gamma\sigma} (\bm{o}_{\alpha})^\sigma
    g_{\delta\lambda} (\bm{o}_{\beta})^\lambda \\
    &= \eta^{\mu\alpha} \eta^{\nu\beta}M^{\gamma\delta}
    g_{\gamma\sigma} (\bm{o}_{\alpha})^\sigma
    g_{\delta\lambda} (\bm{o}_{\beta})^\lambda
\end{align*}
Therefore for the stress-energy tensor we get that
\begin{align*}
    T^{\mu\nu}_{obs}
    &= \eta^{\mu\alpha} \eta^{\nu\beta} T^{\gamma\delta}
    g_{\gamma\sigma} (\bm{o}_{\alpha})^\sigma
    g_{\delta\lambda} (\bm{o}_{\beta})^\lambda
\end{align*}
\item [\textbf{b.}]
Let us consider an observer with four-velocity $\bm{u}_{obs}$ with respect 
to the local rest frame of a perfect fluid.\\
We know that the stress-energy tensor of the perfect fluid is given by equation
(20.15) so in the obeserver's frame the energy density is
\begin{align*}
    T^{tt}_{obs}
    &= \eta^{tt} \eta^{tt} T^{\gamma\delta}
    g_{\gamma\sigma} (\bm{o}_{t})^\sigma
    g_{\delta\lambda} (\bm{o}_{t})^\lambda\\
    &= T^{\gamma\delta}
    g_{\gamma\sigma} (\bm{u}_{obs})^\sigma
    g_{\delta\lambda} (\bm{u}_{obs})^\lambda
\end{align*}
Where we used that $\bm{o}_{t} = \bm{u}_{obs}$, also, we see that in the rest
frame of the perfect fluid $\bm{u}_{obs}$ is
\begin{align*}
    (\bm{u}_{obs})^\mu = \begin{bmatrix}
        \gamma \\ \gamma v_x \\ \gamma v_y \\ \gamma v_z
    \end{bmatrix}
\end{align*}
Also, since the local rest frame is inertial $g_{\mu\nu} = \eta_{\mu\nu}$,
hence
\begin{align*}
    T^{tt}_{obs}
    &= T^{\gamma\delta}
    \eta_{\gamma\sigma} (\bm{u}_{obs})^\sigma
    \eta_{\delta\lambda} (\bm{u}_{obs})^\lambda\\
    &= T^{tt}\gamma^2 + T^{xx}\gamma^2 v_x^2 + T^{yy}\gamma^2 v_y^2
    + T^{zz}\gamma^2 v_z^2\\
    &= \rho_0\gamma^2 + p_0\gamma^2v^2\\
    &= \rho_0\gamma^2 + p_0\gamma^2v^2 + \gamma^2p_0 - \gamma^2p_0\\
    &= \rho_0\gamma^2 - p_0\gamma^2(1 - v^2) + \gamma^2p_0\\
    &= (\rho_0 + p_0)\gamma^2 - p_0
\end{align*}
Finally, if we assume the fluid in the observer's frame moves at the same speed
but in the opposite spatial direction, then the component $u^t$ of the fluid
velocity in the observer's frame is $u^t = \gamma$ and
from equation (20.16) we get that
\begin{align*}
    T^{tt} &= (\rho_0 + p_0)(u^t)^2 + p_0\eta^{tt}
    = (\rho_0 + p_0)\gamma^2 - p_0
\end{align*}
\end{itemize}
\end{proof}

\begin{proof}{\textbf{P20.9}}\\
We know that 
\begin{align*}
    \rho_{obs} = T^{\mu\nu}\eta_{\mu\alpha}\eta_{\nu\beta} u^\alpha_{obs} u^\beta_{obs}
\end{align*}
and that
\begin{align*}
    T^{\mu\nu} = \begin{bmatrix}
        \rho & 0 & 0 & 0\\
        0 & \alpha\rho & 0 & 0\\
        0 & 0 & \alpha\rho & 0\\
        0 & 0 & 0 & \alpha\rho\\
    \end{bmatrix}
    \qquad
    \bm{u}_{obs}^\alpha = \gamma\begin{bmatrix}
        1, v_x, v_y, v_z
    \end{bmatrix}
\end{align*}
Then in the original LIF
\begin{align*}
    \rho_{obs}
    &= T^{tt} (u^t_{obs})^2
    + T^{xx} (u^x_{obs})^2
    + T^{yy} (u^y_{obs})^2
    + T^{zz} (u^z_{obs})^2\\
    &= \rho\gamma^2
    + \alpha\rho\gamma^2v_x^2
    + \alpha\rho\gamma^2v_y^2
    + \alpha\rho\gamma^2v_z^2\\
    &= \rho\gamma^2
    + \alpha\rho\gamma^2v^2\\
    &= \rho\gamma^2
    + \alpha\rho\gamma^2v^2 + \alpha\rho\gamma^2 - \alpha\rho\gamma^2\\
    &= \rho\gamma^2 - \alpha\rho\gamma^2(1-v^2) + \alpha\rho\gamma^2\\
    &= \rho\gamma^2(1 + \alpha) - \alpha\rho
\end{align*}
We want that $\rho(\gamma^2(1 + \alpha) - \alpha) > 0$ then assuming $\rho > 0$
must be that
\begin{align*}
    \gamma^2\alpha - \alpha &> - \gamma^2\\
    \alpha(\gamma^2 - 1) &> -\gamma^2\\
    \alpha &> -\frac{\gamma^2}{\gamma^2 - 1}\\
    \alpha &> \frac{1}{(1- v^2)(1 - \frac{1}{1-v^2})}\\
    \alpha &> \frac{1}{(1- v^2) - 1}\\
    \alpha &> -\frac{1}{v^2}
\end{align*}
Since $\alpha > -1/v^2$ needs to hold for every speed if we set $v = 1$ we get
that
\begin{align*}
    \alpha > -1
\end{align*}
\end{proof}
\end{document}